\documentclass[10pt,a4paper]{amsart}
\usepackage[margin=19mm]{geometry}
\usepackage{amsmath,amssymb,mathtools}
\usepackage[scr=rsfs,cal=euler]{mathalfa}
\usepackage{xfrac}
\usepackage[unicode,hidelinks,bookmarks=false]{hyperref}
\newcommand{\ie}{i.e.\ }
\newcommand{\cf}{cf.\ }
\newcommand{\resp}{respectively\ }
\newcommand{\Subsection}{\subsection}
\providecommand{\nobreakdash}{\nobreak\hbox{-}}
\setlength{\emergencystretch}{2em}
\multlinegap0pt
\usepackage{bm}
\usepackage{bbm}
\usepackage{dsfont}
\usepackage{xspace}
\usepackage{cancel}
\usepackage{mathtools}
\usepackage{amsthm}
\makeatletter
\@ifundefined{thm}{\newtheorem{thm}{Theorem}}{}
\@ifundefined{definition}{\newtheorem{definition}[thm]{Definition}}{}
\makeatother
\newcommand{\dc}{\mathrm{dC}}
\title[Values at non-positive integers of partially twisted multipl]{Values at non-positive integers of partially twisted multiple zeta-functions II}
\author{Driss Essouabri, Kohji Matsumoto, Simon Rutard}
\date{Source: arXiv:2506.20150v1 (CC BY 4.0)}
\begin{document}
\maketitle

\noindent\emph{Source and licence.} Values at non-positive integers of partially twisted multiple zeta-functions II. Version arXiv:2506.20150v1 (CC BY 4.0), distributed under the Creative Commons Attribution 4.0 International licence, as recorded in the arXiv licence metadata for this exact version and shown on the abstract page https://arxiv.org/abs/2506.20150v1. Full rights notice: licenses/arxiv-2506.20150-CC-BY-4.0.txt.\par\smallskip

\noindent\textit{Editorial scope.} Counted unit: the Thm whose source label is \texttt{th:value\_zeta\_n} in the article (source0.tex lines 1061--1070, source label \texttt{th:value\_zeta\_n}), delivered together with the article's own complete original proof of it (source0.tex lines 1350--1402, one \texttt{proof} environment of 53 lines). No mathematics is written, repaired or re-derived: statement and proof are the article's own text line for line. The proof is detached from the statement in the source: the article states the result at lines 1061--1070 and prints its proof at lines 1350--1402, and the proof environment's own header names the statement, so the association is the article's own. Required context retained uncounted: eq:def\_zeta\_de\_crisenoy (lines 288--290); def:F (lines 1011--1019); th:analytic\_continuation\_zeta\_general\_case (lines 1024--1038). Every definition, display and label of those lines is retained unchanged. Omitted from this package and not redistributed: 1--287, 291--1010, 1020--1023, 1039--1060, 1071--1349, 1403--1563. Nothing inside a retained excerpt is omitted or altered: every retained body is delivered exactly as the article's lines are written, so no omission receipt of any kind accompanies this package. Citations: the retained excerpts carry no citation command at all, so no bibliography entry of the article is transcribed and no reference is redistributed. Reference adaptation: none was needed and none was made; every reference inside the delivered unit resolves inside it, so no unresolved reference or citation is produced. Printed slips kept literal: no printed word, symbol, hypothesis or number of the counted statement or of its proof was altered. Layout and class: the article's own class is \texttt{article} with packages including amsmath, amssymb, amsthm, hyperref, textcomp, enumitem, xcolor; the wrapper is the lane's pinned \texttt{amsart} editorial wrapper at 10pt on A4, which restates the theorem-like environments the retained text needs and adds the article's own macro definitions that the retained text needs (\texttt{\textbackslash dc}), with extra wrapper packages (bm, bbm, dsfont, xspace, cancel, mathtools). The delivered numbering is the unit's own. Assets: no retained span contains a figure environment, a graphics-inclusion command, a PDF-page inclusion, a table or a TikZ picture; no image file is copied into this package, the package declares no asset dependency and no image file is redistributed with it. Licence: version arXiv:2506.20150v1 of this article is distributed under the Creative Commons Attribution 4.0 International licence, as recorded in the arXiv licence metadata for this exact version and shown on the abstract page https://arxiv.org/abs/2506.20150v1. A full rights notice is delivered at licenses/arxiv-2506.20150-CC-BY-4.0.txt. Characters: every retained excerpt is pure ASCII, so the delivered engine, XeLaTeX with its default Latin Modern text font, reports no missing character.
\begin{equation}
    \zeta_{n}^{\dc}(\mathbf{s};\mathbf{R};\boldsymbol{\mu}_n) := \sum_{m_1,\ldots,m_n \geq 1} \frac{\mu_1^{m_1} \cdots \mu_n^{m_n}}{R_1(\mathbf{m})^{s_1} \cdots R_T(\mathbf{m})^{s_T}} \label{eq:def_zeta_de_crisenoy}
\end{equation}

\begin{definition}
\label{def:F}
Let $\mathbf{P(Q)}:=(P_1,\ldots,P_{T-1},Q_0,\ldots, Q_d)$. We define the function $F:\mathbb{C}^{T+d} \to \mathbb{C}$ as the holomorphic continuation to the whole complex space of
\begin{align*}
    F(\mathbf{s};\mathbf{z}):= (a_0s_T+\langle \boldsymbol{\alpha},\mathbf{z} \rangle-1) &\zeta(a_0s_T+\langle \boldsymbol{\alpha},\mathbf{z} \rangle) \\
    & \times \zeta_{n-1}^{\dc}\left(s_1,\ldots,s_{T-1},s_T-|\mathbf{z}|,\mathbf{z}; \mathbf{P(Q)};\boldsymbol{\mu}_{n-1}\right),
\end{align*}
where $\mathbf{s}=(s_1,\ldots,s_T) \in \mathbb{C}^T$, $\mathbf{z}=(z_1,\ldots,z_d) \in \mathbb{C}^d$, and where $\zeta^{\dc}$ is defined in \eqref{eq:def_zeta_de_crisenoy}.
\end{definition}

\begin{thm}
\label{th:analytic_continuation_zeta_general_case}
Let $\epsilon>0$ be sufficiently small. Then $\zeta_{n}(\mathbf{s};\mathbf{P};\boldsymbol{\mu}_{n-1})$ has a meromorphic continuation on $\mathcal{H}_{\varepsilon}(\mathbf{M},\boldsymbol{\eta})$ by the form given by
\begin{align}
    \zeta_{n}&(\mathbf{s};\mathbf{P};\boldsymbol{\mu}_{n-1})= \frac{H(\mathbf{s})}{\Gamma(s_T)}+ \sum_{k_1=0}^{M_1} \cdots \sum_{k_d=0}^{M_d} \frac{(-1)^{|\mathbf{k}|}}{\mathbf{k}!} \frac{F(\mathbf{s};-\mathbf{k})\Gamma(s_T+|\mathbf{k}|)}{\Gamma(s_T)(a_0s_T-\langle \boldsymbol{\alpha}|\mathbf{k}\rangle-1)} \label{eq:analytic_continuation_zeta_general_case} \\ 
    & +\sum_{\ell_1=0}^{M_1} \cdots \sum_{\ell_{d-1}=0}^{M_{d-1}} \frac{(-1)^{|\boldsymbol{\ell}|}}{(a_d-a_0)\boldsymbol{\ell}!} \frac{F(\mathbf{s};(-\boldsymbol{\ell},g(s_T;-\boldsymbol{\ell})))\Gamma(\widetilde{g}(s_T;-\boldsymbol{\ell}))\Gamma(g(s_T;-\boldsymbol{\ell}))}{\Gamma(s_T)} \nonumber
\end{align}
with $\mathbf{k}=(k_1,\ldots,k_d)$ and
$\boldsymbol{\ell}=(\ell_1,\ldots,\ell_{d-1})$,
where $H(\mathbf{s})$ is a holomorphic function on $\mathcal{H}_{\varepsilon}(\mathbf{M},\boldsymbol{\eta})$, $F$ is defined in Definition \ref{def:F}, and we set for all $s_T \in \mathbb{C}, \ \mathbf{z}=(z_1,\ldots,z_{d-1}) \in \mathbb{C}^{d-1}$,
\begin{align*}
    g(s_T;\mathbf{z}):=& \frac{1-a_0s_T-\sum_{j=1}^{d-1} (a_j-a_0)z_j}{a_d-a_0}, \\
    \widetilde{g}(s_T;\mathbf{z}):=& s_T-g(s_T;\mathbf{z})-|\mathbf{z}| = \frac{a_ds_T-\sum_{j=1}^{d-1} (a_d-a_j)z_j-1}{a_d-a_0}.
\end{align*}
\end{thm}

\begin{thm}
\label{th:value_zeta_n}
Let $-\mathbf{N}=(-N_1,\ldots,-N_T) \in \mathbb{Z}_{\leq 0}^T$, then
\begin{align}
    \zeta_{n}(-&\mathbf{N};\mathbf{P};\boldsymbol{\mu}_{n-1})=-\sum_{\substack{k_1,\ldots,k_d \geq 0 \\ |\mathbf{k}| \leq N_T}} \binom{N_T}{(\mathbf{k},N_T-|\mathbf{k}|)} \frac{\widetilde{B}_{a_0N_T+\langle \boldsymbol{\alpha}|\mathbf{k}\rangle+1}}{a_0N_T+\langle \boldsymbol{\alpha}|\mathbf{k} \rangle+1} \label{eq:values_zeta_n} \\
    & \hspace{89pt} \times \zeta_{n-1}^{\dc}(-N_1,\ldots,-N_{T-1},-N_T+|\mathbf{k}|,-\mathbf{k};\mathbf{P(Q)};\boldsymbol{\mu}_{n-1}) \nonumber \\
    &+ \frac{N_T!}{a_d}\sum_{i=0}^{\left\lfloor \frac{a_dN_T+1}{\alpha_d}\right\rfloor} \sum_{\substack{\ell_1,\ldots,\ell_{d-1} \geq 0 \\ \sum_{j=1}^{d-1} (a_d-a_j) \ell_j=1+a_dN_T-i\alpha_d}} \frac{(-1)^{|\boldsymbol{\ell}|+i+N_T}(|\boldsymbol{\ell}|+i-N_T-1)!}{i!\boldsymbol{\ell}!} \nonumber \\
    & \hspace{60pt} \times \zeta_{n-1}^{\dc}(-N_1,\ldots,-N_{T-1},-i,-\boldsymbol{\ell},|\boldsymbol{\ell}|+i-N_T;\mathbf{P(Q)};\boldsymbol{\mu}_{n-1}). \nonumber
\end{align}
\end{thm}

\begin{proof}[Proof of Theorem \ref{th:value_zeta_n}]
We apply Theorem \ref{th:analytic_continuation_zeta_general_case} with a tuple $\mathbf{M}=(M_1,\ldots,M_d)\in \mathbb{N}_0^d$ sufficiently large and $\boldsymbol{\eta}=(\eta_1,\ldots,\eta_d) \in (0,1)^d$ such that $-\mathbf{N} \in \mathcal{H}_{\varepsilon}(\mathbf{M},\boldsymbol{\eta})$. We then obtain
\begin{align*}
    \zeta_{n}(\mathbf{s};\mathbf{P};\boldsymbol{\mu}_{n-1})=& -\sum_{k_1=0}^{M_1} \cdots \sum_{k_d=0}^{M_d}  \frac{(-1)^{|\mathbf{k}|}F(-\mathbf{N};-\mathbf{k})}{\mathbf{k}!(a_0N_T+\langle \boldsymbol{\alpha}|\mathbf{k}\rangle+1)}\lim_{s_T \to -N_T}\frac{\Gamma(s_T+|\mathbf{k}|)}{\Gamma(s_T)} \\ 
    & +\sum_{\ell_1=0}^{M_1} \cdots \sum_{\ell_{d-1}=0}^{M_{d-1}} \frac{(-1)^{|\boldsymbol{\ell}|}F(-\mathbf{N};-\boldsymbol{\ell},g(-N_T;-\boldsymbol{\ell}))}{\alpha_d\boldsymbol{\ell}!} \\
    & \hspace{119pt} \times \lim_{s_T \to -N_T}\frac{\Gamma(\widetilde{g}(s_T;-\boldsymbol{\ell}))\Gamma(g(s_T;-\boldsymbol{\ell}))}{\Gamma(s_T)}.
\end{align*} 
It remains to compute the two limits in the formula above.

Let $\mathbf{k}=(k_1,\ldots,k_d) \in \mathbb{N}_0^d$, by the functional relation for $\Gamma(s)$ we get
\[
    \lim_{s_T \to -N_T} \frac{\Gamma(s_T+|\mathbf{k}|)}{\Gamma(s_T)} =
    \begin{cases}
        \frac{(-1)^{|\mathbf{k}|}N_T!}{(N_T-|\mathbf{k}|)!} & \text{ if } |\mathbf{k}| \leq N_T, \\
        0 &\text{ otherwise.}
    \end{cases}
\]

Let $\boldsymbol{\ell}=(\ell_1,\ldots,\ell_{d-1}) \in \mathbb{N}_0^{d-1}$. For the evaluation of the remaining function, we first notice that $\widetilde{g}(-N_T;-\boldsymbol{\ell}) \geq -\frac{a_dN_T+1}{\alpha_d}$. For all $i \in [\![ 0,\lfloor \frac{a_dN_T+1}{\alpha_d} \rfloor ]\!]$, we see that
$\widetilde{g}(-N_T;-\boldsymbol{\ell})=-i$ is equivalent to
\begin{align*}
\sum_{j=1}^{d} (\alpha_d-\alpha_j) \ell_j=1+a_dN_T-i\alpha_d,
\end{align*}
which is further equivalent to
\begin{align*}
    g(-N_T;-\boldsymbol{\ell})=|\boldsymbol{\ell}|+i-N_T.
\end{align*}
It is clear by definition of $g$ that $g(-N_T;-\boldsymbol{\ell}) \in \frac{1}{\alpha_d}\mathbb{N}$, so it is a positive number. Finally, again by the functional relation for $\Gamma(s)$ we get
\[
    \lim_{s_T \to -N_T} \frac{\Gamma(\widetilde{g}(s_T;-\boldsymbol{\ell}))}{\Gamma(s_T)} =
    \begin{cases}
        \frac{(-1)^{N_T}\alpha_dN_T!}{(-1)^ia_di!} & \text{if there exists } 0 \leq i \leq \left\lfloor \frac{a_dN_T+1}{\alpha_d}\right\rfloor \text{ s.t.} \\
        & \sum_{j=1}^{d-1} (a_d-a_j) \ell_j=1+a_dN_T-i\alpha_d, \\
        0 &\text{otherwise.}
    \end{cases}
\]
Therefore, we obtain
\begin{align*}
    \zeta_{n}(&-\mathbf{N};\mathbf{P};\boldsymbol{\mu}_{n-1})=-\sum_{\substack{k_1,\ldots,k_d \geq 0 \\ |\mathbf{k}| \leq N_T}} \binom{N_T}{(k_1,\ldots,k_d,N_T-|\mathbf{k}|)} \frac{F(-\mathbf{N};-\mathbf{k})}{a_0N_T+\langle \boldsymbol{\alpha}|\mathbf{k}\rangle+1} \\
    &+ \sum_{i=0}^{\left\lfloor \frac{a_dN_T+1}{\alpha_d}\right\rfloor} \sum_{\substack{\ell_1,\ldots,\ell_{d-1} \geq 0 \\ \sum_{j=1}^{d-1} (\alpha_d-\alpha_j) \ell_j=1+a_dN_T-i\alpha_d}} \frac{(-1)^{|\boldsymbol{\ell}|+i+N_T}\alpha_dN_T!}{\alpha_da_di!\boldsymbol{\ell}!}(|\boldsymbol{\ell}|+i-N_T-1)! \\
    & \hspace{231pt} \times F(-\mathbf{N};-\boldsymbol{\ell},|\boldsymbol{\ell}|+i-N_T).
\end{align*}
It remains to write explicitly the special values of $F$ appearing in the previous formula. For $\mathbf{k}=(k_1,\ldots,k_d) \in \mathbb{N}_0^d$ such that $|\mathbf{k}|\leq N_T$, by using the relation $\zeta(-m)=-\frac{\widetilde{B}_{m+1}}{m+1}$ we obtain
\begin{align*}
    F(-\mathbf{N};-\mathbf{k})=&\widetilde{B}_{a_0N_T+\langle \boldsymbol{\alpha}|\mathbf{k}\rangle+1} \\
    & \hspace{38pt} \times \zeta_{n-1}^{\dc}(-N_1,\ldots,-N_{T-1},-N_T+|\mathbf{k}|,-\mathbf{k};\mathbf{P(Q)};\boldsymbol{\mu}_{n-1}).
\end{align*}
For $\boldsymbol{\ell}=(\ell_1,\ldots,\ell_{d-1}) \in \mathbb{N}_0^{d-1}$ such that $\alpha_d|\boldsymbol{\ell}|-\langle \boldsymbol{\alpha},\boldsymbol{\ell}\rangle=1+a_dN_T-i\alpha_d$, we have
\begin{align*}
    F(-\mathbf{N};-\boldsymbol{\ell},|\boldsymbol{\ell}&|+i-N_T) \\
    &= \zeta_{n-1}^{\dc}\left((-N_1,\ldots,-N_{T-1},-i,-\boldsymbol{\ell},|\boldsymbol{\ell}|+i-N_T);\mathbf{P(Q)};\boldsymbol{\mu}_{n-1}\right).
\end{align*}
\end{proof}

\end{document}
